% Options for packages loaded elsewhere
\PassOptionsToPackage{unicode}{hyperref}
\PassOptionsToPackage{hyphens}{url}
\documentclass[
]{article}
\usepackage{xcolor}
\usepackage[margin=1in]{geometry}
\usepackage{amsmath,amssymb}
\setcounter{secnumdepth}{-\maxdimen} % remove section numbering
\usepackage{iftex}
\ifPDFTeX
  \usepackage[T1]{fontenc}
  \usepackage[utf8]{inputenc}
  \usepackage{textcomp} % provide euro and other symbols
\else % if luatex or xetex
  \usepackage{unicode-math} % this also loads fontspec
  \defaultfontfeatures{Scale=MatchLowercase}
  \defaultfontfeatures[\rmfamily]{Ligatures=TeX,Scale=1}
\fi
\usepackage{lmodern}
\ifPDFTeX\else
  % xetex/luatex font selection
\fi
% Use upquote if available, for straight quotes in verbatim environments
\IfFileExists{upquote.sty}{\usepackage{upquote}}{}
\IfFileExists{microtype.sty}{% use microtype if available
  \usepackage[]{microtype}
  \UseMicrotypeSet[protrusion]{basicmath} % disable protrusion for tt fonts
}{}
\makeatletter
\@ifundefined{KOMAClassName}{% if non-KOMA class
  \IfFileExists{parskip.sty}{%
    \usepackage{parskip}
  }{% else
    \setlength{\parindent}{0pt}
    \setlength{\parskip}{6pt plus 2pt minus 1pt}}
}{% if KOMA class
  \KOMAoptions{parskip=half}}
\makeatother
\usepackage{color}
\usepackage{fancyvrb}
\newcommand{\VerbBar}{|}
\newcommand{\VERB}{\Verb[commandchars=\\\{\}]}
\DefineVerbatimEnvironment{Highlighting}{Verbatim}{commandchars=\\\{\}}
% Add ',fontsize=\small' for more characters per line
\usepackage{framed}
\definecolor{shadecolor}{RGB}{248,248,248}
\newenvironment{Shaded}{\begin{snugshade}}{\end{snugshade}}
\newcommand{\AlertTok}[1]{\textcolor[rgb]{0.94,0.16,0.16}{#1}}
\newcommand{\AnnotationTok}[1]{\textcolor[rgb]{0.56,0.35,0.01}{\textbf{\textit{#1}}}}
\newcommand{\AttributeTok}[1]{\textcolor[rgb]{0.13,0.29,0.53}{#1}}
\newcommand{\BaseNTok}[1]{\textcolor[rgb]{0.00,0.00,0.81}{#1}}
\newcommand{\BuiltInTok}[1]{#1}
\newcommand{\CharTok}[1]{\textcolor[rgb]{0.31,0.60,0.02}{#1}}
\newcommand{\CommentTok}[1]{\textcolor[rgb]{0.56,0.35,0.01}{\textit{#1}}}
\newcommand{\CommentVarTok}[1]{\textcolor[rgb]{0.56,0.35,0.01}{\textbf{\textit{#1}}}}
\newcommand{\ConstantTok}[1]{\textcolor[rgb]{0.56,0.35,0.01}{#1}}
\newcommand{\ControlFlowTok}[1]{\textcolor[rgb]{0.13,0.29,0.53}{\textbf{#1}}}
\newcommand{\DataTypeTok}[1]{\textcolor[rgb]{0.13,0.29,0.53}{#1}}
\newcommand{\DecValTok}[1]{\textcolor[rgb]{0.00,0.00,0.81}{#1}}
\newcommand{\DocumentationTok}[1]{\textcolor[rgb]{0.56,0.35,0.01}{\textbf{\textit{#1}}}}
\newcommand{\ErrorTok}[1]{\textcolor[rgb]{0.64,0.00,0.00}{\textbf{#1}}}
\newcommand{\ExtensionTok}[1]{#1}
\newcommand{\FloatTok}[1]{\textcolor[rgb]{0.00,0.00,0.81}{#1}}
\newcommand{\FunctionTok}[1]{\textcolor[rgb]{0.13,0.29,0.53}{\textbf{#1}}}
\newcommand{\ImportTok}[1]{#1}
\newcommand{\InformationTok}[1]{\textcolor[rgb]{0.56,0.35,0.01}{\textbf{\textit{#1}}}}
\newcommand{\KeywordTok}[1]{\textcolor[rgb]{0.13,0.29,0.53}{\textbf{#1}}}
\newcommand{\NormalTok}[1]{#1}
\newcommand{\OperatorTok}[1]{\textcolor[rgb]{0.81,0.36,0.00}{\textbf{#1}}}
\newcommand{\OtherTok}[1]{\textcolor[rgb]{0.56,0.35,0.01}{#1}}
\newcommand{\PreprocessorTok}[1]{\textcolor[rgb]{0.56,0.35,0.01}{\textit{#1}}}
\newcommand{\RegionMarkerTok}[1]{#1}
\newcommand{\SpecialCharTok}[1]{\textcolor[rgb]{0.81,0.36,0.00}{\textbf{#1}}}
\newcommand{\SpecialStringTok}[1]{\textcolor[rgb]{0.31,0.60,0.02}{#1}}
\newcommand{\StringTok}[1]{\textcolor[rgb]{0.31,0.60,0.02}{#1}}
\newcommand{\VariableTok}[1]{\textcolor[rgb]{0.00,0.00,0.00}{#1}}
\newcommand{\VerbatimStringTok}[1]{\textcolor[rgb]{0.31,0.60,0.02}{#1}}
\newcommand{\WarningTok}[1]{\textcolor[rgb]{0.56,0.35,0.01}{\textbf{\textit{#1}}}}
\usepackage{graphicx}
\makeatletter
\newsavebox\pandoc@box
\newcommand*\pandocbounded[1]{% scales image to fit in text height/width
  \sbox\pandoc@box{#1}%
  \Gscale@div\@tempa{\textheight}{\dimexpr\ht\pandoc@box+\dp\pandoc@box\relax}%
  \Gscale@div\@tempb{\linewidth}{\wd\pandoc@box}%
  \ifdim\@tempb\p@<\@tempa\p@\let\@tempa\@tempb\fi% select the smaller of both
  \ifdim\@tempa\p@<\p@\scalebox{\@tempa}{\usebox\pandoc@box}%
  \else\usebox{\pandoc@box}%
  \fi%
}
% Set default figure placement to htbp
\def\fps@figure{htbp}
\makeatother
\setlength{\emergencystretch}{3em} % prevent overfull lines
\providecommand{\tightlist}{%
  \setlength{\itemsep}{0pt}\setlength{\parskip}{0pt}}
\usepackage{bookmark}
\IfFileExists{xurl.sty}{\usepackage{xurl}}{} % add URL line breaks if available
\urlstyle{same}
\hypersetup{
  pdftitle={Bayes' Theorem for Parameters},
  hidelinks,
  pdfcreator={LaTeX via pandoc}}

\title{Bayes' Theorem for Parameters}
\author{}
\date{\vspace{-2.5em}2026-04-17}

\begin{document}
\maketitle

\subsection{Introduction}\label{introduction}

Bayesian inference is, mechanically, the application of Bayes' theorem
to a parameter rather than to an event. The prior encodes whatever was
believed before the data arrived, the likelihood is the data-generating
model evaluated at the observed data, and the posterior is the updated
belief. The conceptual leap from frequentist statistics is small but
consequential: parameters are themselves random quantities whose
posterior distribution can be summarised, predicted from, and acted
upon, without invoking long-run sampling thought experiments.

\subsection{Prerequisites}\label{prerequisites}

A working understanding of Bayes' theorem for events, the concept of a
likelihood function, and the difference between probability density and
probability mass.

\subsection{Theory}\label{theory}

For a parameter \(\theta\) and data \(y\), Bayes' theorem reads

\[p(\theta \mid y) = \frac{p(y \mid \theta) p(\theta)}{p(y)} \propto p(y \mid \theta) p(\theta),\]

where the four quantities are the posterior, the likelihood, the prior,
and the marginal likelihood
\(p(y) = \int p(y \mid \theta) p(\theta) \, \mathrm d \theta\). The
marginal likelihood appears as a normalisation constant from a
single-model perspective but plays the central role in
Bayes-factor-based model comparison. For closed-form (conjugate) priors,
the posterior comes out in the same family as the prior with updated
hyperparameters; otherwise MCMC produces draws from the un-normalised
right-hand side.

\subsection{Assumptions}\label{assumptions}

The prior reflects genuine pre-data information or a defensible
weak-regularisation choice, the likelihood correctly captures the
data-generating process, and posterior integrability is established
(proper priors guarantee this; improper priors must be checked).

\subsection{R Implementation}\label{r-implementation}

\begin{Shaded}
\begin{Highlighting}[]
\CommentTok{\# Beta{-}Binomial: prior Beta(2, 2), 8 successes in 10 trials}
\NormalTok{prior\_a }\OtherTok{\textless{}{-}} \DecValTok{2}\NormalTok{; prior\_b }\OtherTok{\textless{}{-}} \DecValTok{2}
\NormalTok{k }\OtherTok{\textless{}{-}} \DecValTok{8}\NormalTok{; n }\OtherTok{\textless{}{-}} \DecValTok{10}
\NormalTok{post\_a }\OtherTok{\textless{}{-}}\NormalTok{ prior\_a }\SpecialCharTok{+}\NormalTok{ k}
\NormalTok{post\_b }\OtherTok{\textless{}{-}}\NormalTok{ prior\_b }\SpecialCharTok{+}\NormalTok{ n }\SpecialCharTok{{-}}\NormalTok{ k}

\NormalTok{theta }\OtherTok{\textless{}{-}} \FunctionTok{seq}\NormalTok{(}\DecValTok{0}\NormalTok{, }\DecValTok{1}\NormalTok{, }\AttributeTok{length.out =} \DecValTok{400}\NormalTok{)}
\FunctionTok{plot}\NormalTok{(theta, }\FunctionTok{dbeta}\NormalTok{(theta, prior\_a, prior\_b), }\AttributeTok{type =} \StringTok{"l"}\NormalTok{,}
     \AttributeTok{main =} \StringTok{"Prior {-}\textgreater{} Posterior"}\NormalTok{, }\AttributeTok{ylab =} \StringTok{"Density"}\NormalTok{, }\AttributeTok{col =} \StringTok{"grey"}\NormalTok{)}
\FunctionTok{lines}\NormalTok{(theta, }\FunctionTok{dbeta}\NormalTok{(theta, k }\SpecialCharTok{+} \DecValTok{1}\NormalTok{, n }\SpecialCharTok{{-}}\NormalTok{ k }\SpecialCharTok{+} \DecValTok{1}\NormalTok{), }\AttributeTok{col =} \StringTok{"orange"}\NormalTok{)  }\CommentTok{\# likelihood scaled}
\FunctionTok{lines}\NormalTok{(theta, }\FunctionTok{dbeta}\NormalTok{(theta, post\_a, post\_b), }\AttributeTok{col =} \StringTok{"steelblue"}\NormalTok{, }\AttributeTok{lwd =} \DecValTok{2}\NormalTok{)}
\FunctionTok{legend}\NormalTok{(}\StringTok{"topleft"}\NormalTok{, }\FunctionTok{c}\NormalTok{(}\StringTok{"Prior"}\NormalTok{, }\StringTok{"Likelihood"}\NormalTok{, }\StringTok{"Posterior"}\NormalTok{),}
       \AttributeTok{col =} \FunctionTok{c}\NormalTok{(}\StringTok{"grey"}\NormalTok{, }\StringTok{"orange"}\NormalTok{, }\StringTok{"steelblue"}\NormalTok{), }\AttributeTok{lty =} \DecValTok{1}\NormalTok{)}

\CommentTok{\# Posterior summaries}
\FunctionTok{c}\NormalTok{(}\AttributeTok{mean =}\NormalTok{ post\_a }\SpecialCharTok{/}\NormalTok{ (post\_a }\SpecialCharTok{+}\NormalTok{ post\_b),}
  \AttributeTok{ci95 =} \FunctionTok{qbeta}\NormalTok{(}\FunctionTok{c}\NormalTok{(}\FloatTok{0.025}\NormalTok{, }\FloatTok{0.975}\NormalTok{), post\_a, post\_b))}
\end{Highlighting}
\end{Shaded}

\subsection{Output \& Results}\label{output-results}

A three-curve plot showing the flat-ish prior, the peaked likelihood,
and the posterior between them, plus a summary giving the posterior mean
(0.71) and the 95 \% equal-tailed credible interval (\(\approx 0.47\) to
\(0.89\)). The posterior is closer to the likelihood than to the prior
because the data (\(n=10\)) carry more precision than the prior
(\(a + b = 4\) pseudo-observations).

\subsection{Interpretation}\label{interpretation}

A reporting sentence: ``Combining a Beta(2, 2) prior with 8 successes in
10 trials yielded a Beta(10, 4) posterior with mean 0.71 (95 \% credible
interval 0.47 to 0.89); the prior contributed the equivalent of about
four pseudo-observations and the data dominated the posterior shape.''
Always disclose the prior; the same data with a different prior would
lead to a slightly different posterior.

\subsection{Practical Tips}\label{practical-tips}

\begin{itemize}
\tightlist
\item
  For conjugate setups (Beta--Binomial, Normal--Normal, Gamma--Poisson,
  Dirichlet--Multinomial), the posterior is closed-form and arithmetic
  alone suffices; reach for MCMC only when conjugacy fails.
\item
  Always report a posterior summary that includes both a central
  tendency and an uncertainty measure; a posterior mean alone is no more
  informative than a frequentist point estimate.
\item
  Sensitivity to the prior matters most in small samples; if your
  conclusion changes with a wider or narrower prior, the data are weak
  and the analysis is genuinely prior-driven.
\item
  Posterior predictive draws (\texttt{rbeta(M,\ post\_a,\ post\_b)}
  followed by \texttt{rbinom(M,\ n\_new,\ theta)}) deliver direct
  uncertainty intervals for future outcomes.
\item
  Visualising prior, likelihood, and posterior on the same plot is the
  single most useful pedagogical artefact when explaining Bayesian
  inference to non-statistical readers.
\end{itemize}

\subsection{R Packages Used}\label{r-packages-used}

Base R for the conjugate update and plotting, \texttt{bayestestR} for
tidy posterior summaries, and \texttt{brms} or \texttt{rstanarm} for
non-conjugate problems where MCMC is required to obtain posterior draws.

\subsection{For Reviewers}\label{for-reviewers}

\textbf{What to look for in a paper using this method.}

\begin{itemize}
\tightlist
\item
  \textbf{Common misapplications.}
\item
  \textbf{Diagnostics that should be reported but often aren't.}
\item
  \textbf{Red flags in tables and figures.}
\item
  \textbf{What to verify.}
\item
  \textbf{What an adequate Methods paragraph must contain.}
\end{itemize}

\subsection{See also --- labs in this
chapter}\label{see-also-labs-in-this-chapter}

\begin{itemize}
\tightlist
\item
  \href{labs/lab-bayesian-thinking-control-vs-decision-error.Rmd}{Bayesian
  thinking; α control vs decision error}
\item
  \href{labs/lab-brms-stan-loo-hierarchical-models.Rmd}{brms / Stan,
  LOO, hierarchical models}
\end{itemize}

\end{document}
